\begin{lemma}[6.15.8]
\label{IV.2.s6.15.8-fr}
Sean $k$ un cuerpo separablemente cerrado \textup{(es decir, tal
que la clausura algebraica de $k$ sea \emph{radicial} sobre $k$)},
$X$ un $k$-preesquema localmente de tipo finito sobre $k$ y
$x$ un punto \emph{cerrado} de $X$. Si $X$ es unirramoso en el punto
$x$, es geométricamente unirramoso en ese punto.
\end{lemma}

\begin{proof}
En efecto, se sabe (\textbf I, 6.4.2) que $\boldsymbol{k}(x)$ es una extensión
\emph{algebraica} de $k$; como el cuerpo residual de la clausura
integral de $(\mathcal O_x)_{\mathrm{red}}$ (clausura
integral que, por hipótesis, es un anillo local) es una
extensión algebraica de $\boldsymbol{k}(x)$, es una extensión \emph{radicial}
de $k$ por hipótesis y, por tanto, también de $\boldsymbol{k}(x)$.
\end{proof}
